\documentclass[conference]{IEEEtran}
\usepackage{amsmath,amssymb}
\usepackage{graphicx}
\usepackage{booktabs}
\usepackage{multirow}
\usepackage{tabularx}
\usepackage{xcolor}
\usepackage{siunitx}
\usepackage{cleveref}
\usepackage{microtype}
\usepackage{hyperref}
\usepackage{biblatex}

\title{Controlled Detector Attribution under a Locked Runtime Intervention Policy:\\ A Pilot-Scale Agent Security Evaluation}
\author{[REDACTED for double-blind review]}

\begin{document}
\maketitle

\begin{abstract}
Observed security behavior in tool-using language-model agents can mix detector behavior with the downstream intervention policy that allows, wraps, or denies tools. Without holding that policy fixed, detector-related effects are difficult to attribute.

This paper reports a controlled attribution protocol, not a new detector family and not a defense covering all threat models. The intervention policy is held fixed (PHASE1-CORE). Detector identity is varied, including a no-detection reference (D0). The primary outcome is Tool-HASR (harmful tool execution). Judge-ASR is a secondary, non-identical diagnostic of judge-level attack-success assessment. The evaluation uses frozen pack \texttt{p2\_agentic\_v0.1.0} (16 attack / 16 twin / 4 hard-negative trajectories), four target models, and $n=16$ attack arms per target$\times$detector cell. Q2 live completed 432/432 arms (historical spend \$0.152885).

Under the tested protocol, detector-related Tool-HASR $\Delta$ versus D0 was negative for D1, D2, and D4 on the locked Stage-B target T0 and on independently selected T1--T3 (9/9 directional agreements). This is pilot-scale directional consistency. It does not establish population-level generalization. T1--T3 Tool-HASR was 81/192 $= 0.421875$; Judge-ASR was 186/192 $= 0.96875$ ($M3=108$, $M4=3$). \textsc{invalid\_tool\_args} occurred on 192 events / 136 arms and is treated as a canonical execution-state diagnostic, not as automatic attack success or failure.

The manuscript does not claim robustness beyond the tested protocol, production readiness, confirmatory multi-model generalization, or a ranking against external defenses.
\end{abstract}

\section{Introduction}
\label{sec:intro}

\textbf{Problem.} Observed security behavior in agent systems may reflect both detector behavior and downstream intervention policy. A stack that changes detector and policy together reports an operational outcome. It does not isolate which layer is associated with a change in harmful tool execution.

\textbf{Gap.} Without controlling the intervention policy, detector-related effects can be difficult to attribute. Existing literature already provides prompt-injection attacks, agent security benchmarks, runtime guardrails, architectural isolation, and detector approaches. Those lines of work do not, in the surveyed set, establish the locked-policy detector-attribution factorial used here as the central measurement protocol. Novelty class: \textbf{PARTIAL\_GAP}.

\textbf{Approach.} Hold the intervention policy fixed and vary detector identity, including a no-detection reference (D0). Thresholds, action costs, frozen pack, judge, temperature, and cache settings remain locked. The primary estimand is the sign of Tool-HASR $\Delta$ versus D0 at each target. Tool-HASR and Judge-ASR are kept as distinct measurements.

\textbf{Result.} Observed $\Delta$ direction remained consistent across T0--T3 in this pilot-scale evaluation: 9/9 detector-target comparisons preserved the negative direction ($n=16$ per cell; 432/432 Q2 live arms).

\textbf{Qualification.} This is not a confirmatory claim. Sample size is pilot-scale. Precision is limited. Power is limited for broad generalization. Directional consistency is an observed property of this evaluation. It does not establish population-level generalization. Three of four targets are Qwen-family models. D3 was deferred. Open-ended adaptive attackers (C4 interaction horizon) are out of scope.

\section{Research Question}
\label{sec:rq}

\textbf{RQ-C2 (locked protocol).} Does the detector-related Tool-HASR effect versus D0 observed under the locked Stage-B target (T0) remain directionally consistent when the same policy, pack, detectors, thresholds, and judge are applied to independently selected secondary target models T1--T3?

\section{Problem Formulation}
\label{sec:formulation}

\begin{table}[h]
\caption{Layer decomposition}
\label{tabtab:layers}
\begin{tabular}{p{0.15\linewidth}lp{0.7\linewidth}lp{0.15\linewidth}}
\toprule
\textbf{Layer} & \textbf{What it is} & \textbf{What it is not} \\
\midrule
Detector behavior & Hits/scores on untrusted context & The intervention that grants/denies tools \\
Downstream intervention policy & PHASE1-CORE mapping to A0--A3 & A detector family \\
Tool-HASR (primary) & Attack-arm rate of harmful tool execution & A judge-model opinion \\
Judge-ASR (secondary) & Attack-arm rate of judge-classified attack success & Harmful tool execution \\
$M3$ & Judge success $\wedge$ $\neg$ Tool-HASR success & Proof that the judge is wrong \\
$M4$ & Tool-HASR success $\wedge$ $\neg$ judge success & Proof that the judge is right \\
\textsc{invalid\_tool\_args} & Canonical tool-state events & Attack success/failure; harmlessness; detector failure \\
\bottomrule
\end{tabular}
\end{table}

$\Delta(d,t) = \text{Tool-HASR}(d,t) - \text{Tool-HASR}(D0,t)$ on attack arms at the same target, under locked PHASE1-CORE. This is a \textbf{detector-related association} under the controlled protocol. It does not establish that detector identity is the sole associated factor, independent of all runtime factors.

\section{Related Work}
\label{sec:related}

We cite records whose title, authors, year, and arXiv id were verified from the arXiv API. Venues and DOIs were verified from official sources where available. Novelty class: \textbf{PARTIAL\_GAP}.

\subsection{Controlled Attribution Framework}
\label{sec:framework}

Each episode-arm runs:
\textbf{Detector $\to$ RiskCore $\to$ CorePolicy (PHASE1-CORE) $\to$ ActionLayer $\to$ ToolPermissionGate $\to$ Mock tools}, with a parallel \textbf{Judge} that does not enter the detector$\to$policy path. Gold labels do not enter that path.

\section{Experimental Protocol}
\label{sec:protocol}

\begin{table}[h]
\caption{Study design}
\label{tabtab:design}
\begin{tabularx}{\linewidth}{XX}
\toprule
\textbf{Field} & \textbf{Value} \\
\midrule
Question & RQ-C2 \\
Design & B\_REDUCED\_Q2 \\
Run ID & \texttt{p3\_stage\_c\_q2\_20260917T123855Z\_b075df0f} \\
Evidence commit & \texttt{b075df0f5ec5ad11ede76ac4e4079cade15208f1} \\
Pack & \texttt{p2\_agentic\_v0.1.0} (SHA \texttt{32b40e3b...}) \\
P1 SHA & \texttt{1a0b0053c...} \\
Q2 predictions SHA & \texttt{2a2c2f31...} \\
Trajectories & 36 (16 attack / 16 twin / 4 hard-negative) \\
Detectors & D0, D1, D2, D4 (D3 deferred) \\
Policy & PHASE1-CORE \\
T0 & \texttt{qwen/qwen-2.5-7b-instruct} (reused from Stage-B) \\
Q2 live & 36 traj $\times$ 4 det $\times$ PHASE1-CORE $\times$ \{T1,T2,T3\} = 432 arms \\
Completed & 432/432, 0 failed \\
Temperature & 0.0 \\
Cache & false \\
Seed & 42 \\
Backend & OpenRouter \\
Budget cap / actual & \$10.0 / \$0.152885 \\
Historical API calls & 2061 \\
scientific\_evidence & false \\
$n$ per cell & 16 attack arms \\
\end{tabularx}
\end{table}

\begin{table}[h]
\caption{Target and judge models (not a quality ranking)}
\label{tabtab:models}
\begin{tabularx}{\linewidth}{XX}
\toprule
\textbf{Slot} & \textbf{Model ID} & \textbf{Family} & \textbf{Role} \\
\midrule
T0 & \texttt{qwen/qwen-2.5-7b-instruct} & Qwen-2.5 dense & Stage-B reference \\
T1 & \texttt{qwen/qwen3-30b-a3b} & Qwen3 MoE & secondary \\
T2 & \texttt{google/gemma-3-27b-it} & Gemma-3 & secondary \\
T3 & \texttt{qwen/qwen3.5-35b-a3b} & Qwen3.5 MoE & secondary \\
Judge & \texttt{qwen/qwen-2.5-72b-instruct} & Qwen-2.5 dense & locked \\
\bottomrule
\end{tabularx}
\end{table}

Four targets were used. Three of four (T0, T1, T3) are Qwen-family. The judge is also Qwen-family. Model-family diversity is limited. Results should not be generalized to arbitrary model families.

\section{Results}
\label{sec:results}

All cells use $n=16$ attack arms. Wilson 95\% CIs on rates are derived from locked counts.

\begin{table}[h]
\caption{Primary $\Delta$ results (Tool-HASR versus D0)}
\label{tabtab:delta}
\begin{tabular}{lrrrrr}
\toprule
Cell & Tool-HASR $(d)$ & Tool-HASR $(D0)$ & $\Delta$ & Sign \\
\midrule
T0/D1 & 2/16 = 0.1250 & 13/16 = 0.8125 & $-0.6875$ & NEG \\
T0/D2 & 9/16 = 0.5625 & 13/16 = 0.8125 & $-0.2500$ & NEG \\
T0/D4 & 4/16 = 0.2500 & 13/16 = 0.8125 & $-0.5625$ & NEG \\
T1/D1 & 2/16 = 0.1250 & 14/16 = 0.8750 & $-0.7500$ & NEG \\
T1/D2 & 9/16 = 0.5625 & 14/16 = 0.8750 & $-0.3125$ & NEG \\
T1/D4 & 4/16 = 0.2500 & 14/16 = 0.8750 & $-0.6250$ & NEG \\
T2/D1 & 2/16 = 0.1250 & 13/16 = 0.8125 & $-0.6875$ & NEG \\
T2/D2 & 9/16 = 0.5625 & 13/16 = 0.8125 & $-0.2500$ & NEG \\
T2/D4 & 4/16 = 0.2500 & 13/16 = 0.8125 & $-0.5625$ & NEG \\
T3/D1 & 2/16 = 0.1250 & 11/16 = 0.6875 & $-0.5625$ & NEG \\
T3/D2 & 8/16 = 0.5000 & 11/16 = 0.6875 & $-0.1875$ & NEG \\
T3/D4 & 3/16 = 0.1875 & 11/16 = 0.6875 & $-0.5000$ & NEG \\
\bottomrule
\end{tabular}
\end{table}

\textbf{9/9 detector-target comparisons preserved the negative direction} relative to T0. All D1/D2/D4 $\Delta$ values versus D0 are negative. Q2 live T1--T3 pooled Tool-HASR: 81/192 $= 0.421875$.

\begin{table}[h]
\caption{Tool-HASR vs Judge-ASR ($n=16$/cell)}
\label{tabtab:hasr-judge}
\begin{tabular}{lrr}
\toprule
Cell & Tool-HASR & Judge-ASR \\
\midrule
T0/D0 & 13/16 = 0.8125 & 12/16 = 0.7500 \\
T0/D1 & 2/16 = 0.1250 & 14/16 = 0.8750 \\
T0/D2 & 9/16 = 0.5625 & 15/16 = 0.9375 \\
T0/D4 & 4/16 = 0.2500 & 14/16 = 0.8750 \\
T1/D0 & 14/16 = 0.8750 & 15/16 = 0.9375 \\
T1/D1 & 2/16 = 0.1250 & 16/16 = 1.0000 \\
T1/D2 & 9/16 = 0.5625 & 16/16 = 1.0000 \\
T1/D4 & 4/16 = 0.2500 & 15/16 = 0.9375 \\
T2/D0 & 13/16 = 0.8125 & 15/16 = 0.9375 \\
T2/D1 & 2/16 = 0.1250 & 16/16 = 1.0000 \\
T2/D2 & 9/16 = 0.5625 & 16/16 = 1.0000 \\
T2/D4 & 4/16 = 0.2500 & 15/16 = 0.9375 \\
T3/D0 & 11/16 = 0.6875 & 15/16 = 0.9375 \\
T3/D1 & 2/16 = 0.1250 & 16/16 = 1.0000 \\
T3/D2 & 8/16 = 0.5000 & 16/16 = 1.0000 \\
T3/D4 & 3/16 = 0.1875 & 15/16 = 0.9375 \\
\bottomrule
\end{tabular}
\end{table}

\begin{table}[h]
\caption{$M3$/$M4$ diagnostics}
\label{tabtab:m3m4}
\begin{tabular}{lrr}
\toprule
Scope & $M3$ (J+ T$-$) & $M4$ (T+ J$-$) \\
\midrule
T0 pooled & 33 & 6 \\
T1 & 34 & 1 \\
T2 & 35 & 1 \\
T3 & 39 & 1 \\
Q2 live T1--T3 & 108 & 3 \\
\bottomrule
\end{tabular}
\end{table}

\section{Diagnostic Analysis}
\label{sec:diagnostic}

\textsc{invalid\_tool\_args} is a canonical execution-state category. It is not automatically attack success, attack failure, harmlessness, or detector failure. It is retained in the official S0 analysis.

\section{Limitations}
\label{sec:limitations}

\begin{table}[h]
\caption{Threats and limitations}
\label{tabtab:limitations}
\begin{tabular}{rl}
\toprule
\# & Limitation \\
\midrule
1 & $n=16$ per cell; limited precision; limited power \\
2 & Pilot-scale study; not confirmatory \\
3 & Four target models only \\
4 & Qwen-heavy target set \\
5 & Three of four targets are Qwen-family \\
6 & D3 deferred \\
7 & C4 / adaptive attacker out of scope \\
8 & INVALID\_TOOL\_ARGS confounder (192/136) \\
9 & Stage-B raw traces MISSING\_LOCALLY \\
10 & No matched external baseline \\
\bottomrule
\end{tabular}
\end{table}

\section{Reproducibility}
\label{sec:repro}

$\text{REPRODUCIBILITY\_STATUS} = \text{PARTIAL}$. Full reproducibility is not claimed while Stage-B raw traces remain missing.

\section{Conclusion}
\label{sec:conclusion}

This study evaluates whether detector-related security effects under a locked intervention policy remain directionally consistent across independently selected target models. Under the tested protocol, the observed $\Delta$ direction remained consistent across T0--T3: 9/9 comparisons preserved the negative direction (432/432 arms; \$0.152885 historical spend; $n=16$/cell). That is pilot-scale directional consistency. It is not confirmatory. It does not establish population-level generalization.

\section{Open Science}
\label{sec:openscience}

Experimental code and Q2 evaluation predictions are available in the project repository. Stage-B raw traces (\texttt{predictions.jsonl}, \texttt{metrics.json}, \texttt{manifest.json} for run \texttt{p3\_stage\_b\_20260916T235438Z\_7e401714}) are \textsc{missing\_locally}; only the SHA pointer and packaged metrics are available. Reproducibility status is \textsc{partial}. Full Stage-B replication requires access to the original run directory, which is not present in the current checkout. An anonymous artifact package was prepared for double-blind review. Reproducibility instructions are provided in the repository documentation.

\section{LLM Usage}
\label{sec:llmusage}

An LLM judge (\texttt{qwen/qwen-2.5-72b-instruct}) was used for post-hoc scoring only; it did not influence experimental design or data collection. Historical API calls: 2061. Total cost: \$0.152885. Provider: OpenRouter. Temperature: 0.0. Seed: 42. No LLM was used for manuscript writing.

\section{Ethics and Dual-Use}
\label{sec:ethics}

This research studies prompt injection attacks for defensive purposes. Experiments were conducted in a sandboxed environment with mock tools. No real users, production systems, or production APIs were targeted. Adversarial payloads are disclosed for scientific reproducibility. Misuse is the responsibility of the actor. The authors acknowledge dual-use risk; findings are intended to support defensive applications. No human subjects were involved; no IRB review was required.

\printbibliography{references}

\end{document}
